\section{EqRel (Relações de Equivalência)}

As construções em \textbf{EqRel} lidam com a manipulação de classes de equivalência para formar novos conjuntos quocientes ou uniões estruturadas.

\subsection{Fundamentação Teórica}
Nesta categoria, destacam-se as seguintes construções:
\begin{itemize}
    \item \textbf{Produtos:} O produto de duas relações $R_A$ e $R_B$ relaciona pares de elementos se as suas respetivas componentes estiverem relacionadas nas estruturas originais.
    \item \textbf{Coprodutos:} Definidos pela união disjunta das relações, onde não existe relação cruzada entre elementos de conjuntos diferentes.
\end{itemize}

\subsection{Implementação Computacional}
A implementação do produto baseia-se na interseção das relações estendidas ao novo espaço cartesiano.

\begin{lstlisting}[caption={Produto de Relações de Equivalência}]
function R_prod = eqrel_product(RA, RB)
    % RA e RB são matrizes simétricas de equivalência
    % O produto preserva a estrutura de blocos das classes
    R_prod = kron(RA, RB); 
end
\end{lstlisting}